%!TEX program=xelatex
\documentclass{ctexart}

\usepackage{amscd,amsthm,amsfonts,amssymb,amsmath}
\usepackage[colorlinks=true]{hyperref}
\usepackage[all]{xy}
\usepackage{mathrsfs}
\usepackage{tikz,mathpazo}
\usepackage{bm}
\usepackage{geometry}
\usepackage{xcolor}
\usepackage{eso-pic}
%\usepackage{CJK}
\usepackage{arydshln}
\usepackage{mathptmx}
\usepackage[11pt]{moresize}

\newtheorem{thm}{定理}
\newtheorem{cor}[thm]{推论}
\newtheorem{lem}[thm]{引理}
\newtheorem{prop}[thm]{命题}
\newtheorem{defn}[thm]{定义}
\newtheorem{ques}[thm]{问题}
\newtheorem{eg}[thm]{例题}
\newtheorem{ex}[thm]{习题}
\newtheorem{rmk}[thm]{注释}
\newtheorem{ntn}[thm]{记号}

\newtheorem{exercise}[thm]{习题}

\allowdisplaybreaks
\raggedbottom

\usepackage{mathdots}
\usepackage{extarrows}
\usepackage{amsbsy}
\newcommand\bs[1]{\boldsymbol{ #1}}
\newcommand\hint[1]{\par {\footnotesize\kaishu 提示：#1}}
\newcommand\note[1]{\par {\footnotesize\kaishu 注意：#1}}
		
\def\hard{$\heartsuit$}
\def\zero{O}
\def\dps{\displaystyle}
\def\wtd{\widetilde}
\def\what{\widehat}
\DeclareMathOperator{\T}{T}
\DeclareMathOperator{\trace}{trace}
\DeclareMathOperator{\rank}{rank}
\DeclareMathOperator{\toep}{toep}
\DeclareMathOperator{\range}{\mathcal{R}}
\DeclareMathOperator{\nullspace}{\mathcal{N}}
\newcommand{\diff}{\mathop{}\!\mathrm{d}}
\DeclareMathOperator{\subspan}{span}

\usepackage{mathtools}
\providecommand\given{\:\vert\:}
\DeclarePairedDelimiterXPP\set[1]{}{\{}{\}}{}{
\renewcommand\given{\nonscript\:\delimsize\vert\nonscript\:\mathopen{}}
#1}
\DeclarePairedDelimiter\abs{|}{|}
\DeclarePairedDelimiter\N{\|}{\|}
\DeclarePairedDelimiter\spanbasic{(}{)}
\newcommand\lspan[1]{\subspan\spanbasic*{#1}}

\newif\ifhideproofs
\hideproofstrue
\newif\iffinalver
%\finalvertrue        
\ifhideproofs
    \usepackage[final]{changes}
    \definechangesauthor[name=Error, color=red]{err}
    \usepackage{environ}
    \NewEnviron{hide}{}
    \let\proof\hide
    \let\endproof\endhide
\else
    \iffinalver
        \usepackage[final]{changes}
    \else
        \usepackage{changes}
    \fi
    \definechangesauthor[name=Error, color=red]{err}
\fi

\begin{document}


\title{习题课材料（九）}
\author{}
\date{}
\maketitle

\begin{exercise}
设$A, B$是$n$阶方阵, 且满足$AB=BA$. 
\begin{enumerate}
\item 证明, 若$A$有$n$个不同的特征值, 则$B$可对角化. 
\item 若$A$有代数重数大于$1$的特征值, $B$是否一定可对角化？
\end{enumerate}
\end{exercise}
\begin{proof}
	\begin{enumerate}
		\item $A$有$n$个不同特征值, 则$A$可对角化, 令$A=X\Lambda X^{-1}$. 
			则$AB=BA\Rightarrow \Lambda X^{-1}BX=X^{-1}BX\Lambda$. 
			记$\wtd B=X^{-1}BX$, 
			$\Lambda \wtd B=\wtd B\Lambda$ 说明$\wtd B$是对角矩阵. 
			（事实上, 逐个元素计算, 有$(\lambda_i-\lambda_j)\wtd B_{ij}=0\Rightarrow \wtd B_{ij}=0,i\ne j$. ）
			因此$B=X\wtd BX^{-1}$可对角化. 
			\item 不一定. 例如$A=I$, 则$B$可以是任意矩阵, 包括不可对角化的矩阵. 
	\end{enumerate}
\end{proof}

\begin{exercise} \hard
给定$n$阶方阵$A, B$满足$AB=BA$. 
\begin{enumerate}
\item 证明, 若$A$有$n$个不同的特征值, 则存在次数不超过$n-1$的多项式$f(x)$, 使得$B=f(A)$. 
\item 证明, 若$A=J_n(\lambda)$, 则存在次数不超过$n-1$的多项式$f(x)$, 使得$B=f(A)$. 
\item  举例说明, 存在$A, B$满足$AB=BA$, 但不存在多项式$f(x)$, 使得$B=f(A)$. 
\end{enumerate}
\end{exercise}
\begin{proof}
	\begin{enumerate}
		\item 利用上题, $A,B$可以同时对角化, 令$A=X\Lambda X^{-1},B=XMX^{-1}$. 
			则$B=f(A)\Leftrightarrow M=f(\Lambda) \Leftrightarrow \mu_i=f(\lambda_i),i=1,\dots,n$. 这是例4.2.13讲过的插值问题, 解存在唯一. 
		\item $A=J_n(\lambda)=\lambda I +J_n(0)$, 因此$J_n(0)B=BJ_n(0)$, 即
			\begin{align*}
				\begin{bmatrix}
					b_{21} & b_{22} & \cdots & b_{2n}\\
					\vdots & \vdots &  & \vdots \\
					b_{n1} & b_{n2} & \cdots & b_{nn}\\
					\\
				\end{bmatrix}
				&=\begin{bmatrix}
					& 1\\ && \ddots \\ &&& 1\\ \\
				\end{bmatrix}\begin{bmatrix}
					b_{11} & b_{12} & \cdots & b_{1n}\\
					b_{21} & b_{22} & \cdots & b_{2n}\\
					\vdots & \vdots &  & \vdots \\
					b_{n1} & b_{n2} & \cdots & b_{nn}\\
				\end{bmatrix}
				\\&=\begin{bmatrix}
					b_{11} & b_{12} & \cdots & b_{1n}\\
					b_{21} & b_{22} & \cdots & b_{2n}\\
					\vdots & \vdots &  & \vdots \\
					b_{n1} & b_{n2} & \cdots & b_{nn}\\
				\end{bmatrix}\begin{bmatrix}
					& 1\\ && \ddots \\ &&& 1\\ \\
				\end{bmatrix}
				={\color{blue}\begin{bmatrix}
                    & b_{12} & b_{13} & \cdots \\
                    & 	b_{22} & b_{23} & \cdots \\
                    & 	\vdots & \vdots &   \\
                    & 	b_{n2} & b_{n3} & \cdots \\
                    \end{bmatrix}}
			\end{align*}
            \comment[id=err]{以上蓝色部分原先元素下标有误}
			比较得到
			\[
				B=\begin{bmatrix}
				b_{11} & b_{12} & \cdots & b_{1n}\\
				 & b_{11} & \cdots & b_{1,n-1}\\
				 &  & \ddots & \vdots \\
				 &  &  & b_{11}\\
			 \end{bmatrix}=b_{11}I+b_{12}J_n(0)+\dots+b_{1n}J_n(0)^{n-1}
			 =b_{11}I+b_{12}(A-\lambda I)+\dots +b_{1n}(A-\lambda I)^{n-1}
			 .
			\]
		 \item 类似上题, 取$A=I,B=J_n(0)$, 则二者可交换, 但对任意多项式$f$, $f(A)$是对角矩阵, 不可能等于$B$. 
	\end{enumerate}
\end{proof}

\begin{exercise}
\hard\hard 证明, 任意迹为$0$的方阵相似于一个对角元素全为$0$的方阵. 

%\hint{利用数学归纳法. 
%先取非零向量$\bs q$且$\bs q$不是$A$的特征向量, 再把$\bs q,A\bs q$扩充成一组基做相似变换. }
\end{exercise}
\begin{proof}
	记迹为$0$的方阵为$A$. 
	阶数为$1$时显然. 用数学归纳法. 目标是求$Q$使得$AQ=Q\begin{bmatrix}
		0 & *\\ * & \wtd A
	\end{bmatrix}$. 
	展开计算, 先观察$Q=\begin{bmatrix}
		\bs q_1 & Q_2
	\end{bmatrix}$, 可知$A\bs q_1=Q_2*$. 也就是说, $A\bs q_1$可用$Q_2$的列线性表示, 而不含$\bs q_1$. 
	这提示我们如下构造$Q$：取一个向量$\bs q$, 它不是$A$的特征向量, 则$\bs q,A\bs q$线性无关. 将之扩充为空间的一组基, 按列组成矩阵$Q$, 就有$AQ=Q\begin{bmatrix}
		0 &*\\ \bs e_1 & \wtd A
	\end{bmatrix}$. 
	利用归纳假设, 存在$P$使得$\wtd AP=P\what A$, 其中$\what A$的对角元素全是$0$, 那么
	$AQ\begin{bmatrix}
		1 & \\ & P
	\end{bmatrix}=Q\begin{bmatrix}
	1 & \\ & P
	\end{bmatrix}\begin{bmatrix}
	0 & *\\ P^{-1}\bs e_1 &\what A
	\end{bmatrix}$, 也满足条件. 
	另一方面, 如果不存在这样的$\bs q$, 那么任意向量都是$A$的特征向量. 
	$A\bs e_i=\lambda_i\bs e_i$, 即$A$是对角阵. 
	再看$A(\bs e_1+\dots+\bs e_n)=\lambda(\bs e_1+\dots+\bs e_n)$, 可知$\lambda_1=\dots=\lambda_n=\lambda$, 即$A=\lambda I$. 
	但$\trace(A)=n\lambda=0$, 于是$A=0$. 
\end{proof}

\begin{exercise}\label{excs:similar:transpose}
利用Jordan标准形证明, $A$与$A^{\T}$相似. 
\end{exercise}
\begin{proof}
	设$A=XJX^{-1}$, 其中$X$可逆, $J$是$A$的Jordan标准形. 
	则$A^T=X^{-\T}J^TX^T$. 只要能证明对任意Jordan块$J_n(\lambda)$, $J_n(\lambda)$与$J_n(\lambda)^T$相似, 就有$J$与$J^T$相似, 从而$A$与$A^T$相似. 

	下证对Jordan块$J_n(\lambda)$命题成立. 事实上证明对特殊Jordan块$J_n(0)$成立即可. 
	取$F=\begin{bmatrix}
		&& 1\\ & \iddots\\ 1\\
	\end{bmatrix}$, 不难验证, 
	\[
		FJ_n(0)F^{-1}=\begin{bmatrix}
		&& 1\\ & \iddots\\ 1\\
	\end{bmatrix}\begin{bmatrix}
				& 1\\ && \ddots \\ &&& 1\\ \\
	\end{bmatrix}\begin{bmatrix}
		&& 1\\ & \iddots\\ 1\\
	\end{bmatrix}=\begin{bmatrix}
				\\ 1\\ & \ddots \\ && 1&\\ 
			\end{bmatrix}=J_n(0)^T.
	\]
\end{proof}

\begin{exercise}\label{exc:jordan-simplify}
给定$m$阶方阵$A_1$, $n$阶方阵$A_2$和$m\times n$矩阵$B$. 
证明：
\begin{enumerate}
\item 如果$A_1$和$A_2$没有相同的特征值, 则关于$m\times n$矩阵$X$的Sylvester方程$A_1X-XA_2=B$有唯一解. 
\item 如果$A_1$和$A_2$没有相同的特征值, 则存在唯一的矩阵$X$满足
\[
\begin{bmatrix}
I_m & X\\ \zero & I_n\\
\end{bmatrix}\begin{bmatrix}
A_1 & B\\ \zero &A_2\\
\end{bmatrix}\begin{bmatrix}
I_m & X\\ \zero & I_n\\
\end{bmatrix}^{-1}=\begin{bmatrix}
A_1 & \zero\\ \zero &A_2\\
\end{bmatrix}.
\]
\item \label{itm:exc:jordan-simplify:3}
	对$n$阶方阵$A$, 存在可逆矩阵$X$, 使得
\[
X^{-1}AX=\begin{bmatrix}
	\lambda_1I+N_1 &&        &                \\
&\lambda_1I+N_2 &        &                \\
&               & \ddots &                \\
&               &        & \lambda_sI+N_s \\
\end{bmatrix},
\]
其中$N_1,N_2,\dots,N_s$是严格上三角矩阵. 
\item \hard 如果$m>n$且$A_1=J_m(0),A_2=J_n(0),B=\bs e_1\bs b^{\T}$, 即$B$除第一行外元素全为$0$, 则关于$m\times n$矩阵$X$的Sylvester方程$A_1X-XA_2=B$有解, 于是存在矩阵$X$满足
\[
\begin{bmatrix}
I_m & X\\ \zero & I_n\\
\end{bmatrix}\begin{bmatrix}
A_1 & B\\ \zero &A_2\\
\end{bmatrix}\begin{bmatrix}
I_m & X\\ \zero & I_n\\
\end{bmatrix}^{-1}=\begin{bmatrix}
A_1 & \zero\\ \zero &A_2\\
\end{bmatrix}.
\]
\item \label{itm:exc:jordan-simplify:5}
对$n$阶方阵$A=\begin{bmatrix}
		J_{k_1}(0) & \bs e_1\bs a_2^{\T} & \cdots & \bs  e_1\bs a_r^{\T}\\
		&J_{k_2}(0) \\
		&& \ddots \\
		&&&J_{k_r}(0)\\
	\end{bmatrix}$, 其中$k_1> k_2\ge\dots\ge k_r$（即$A$是一个Jordan标准形与除第一行外元素全为$0$的矩阵的和）, 
存在可逆矩阵$X$, 使得
\[
X^{-1}AX=\begin{bmatrix}
		J_{k_1}(0) \\
		&J_{k_2}(0)  \\
		&& \ddots \\
		&&&J_{k_r}(0)\\
\end{bmatrix}.
\]
\end{enumerate}
\end{exercise}
\begin{proof}
	\begin{enumerate}
		\item 前面习题已经证明过$A_2$上三角的特殊情形. 注意到, 对任意矩阵$A_2$, 存在可逆矩阵$Y$, 使得$Y^{-1}A_2Y=T$是上三角矩阵. 因此$A_1X-XA_2=B\Leftrightarrow A_1X-XYTY^{-1}=B\Leftrightarrow (Y^{-1}A_1Y)(Y^{-1}XY)-(Y^{-1}XY)T=Y^{-1}BY$, 就转化为已知情形. 
			\item $\begin{bmatrix}
					A_1 & \\ & A_2
			\end{bmatrix}=\begin{bmatrix}
			I & X\\& I
			\end{bmatrix}\begin{bmatrix}
			A_1 & B\\ & A_2
			\end{bmatrix}\begin{bmatrix}
			I & -X \\ & I
			\end{bmatrix}=\begin{bmatrix}
			A_1 & B+XA_2-A_1X\\ & A_2
			\end{bmatrix}\Leftrightarrow B+XA_2-A_1X=0
			$, 就化为上一小题. 
		\item 先利用可逆矩阵$X_1$把$A$上三角化, 即$X_1^{-1}AX_1=\begin{bmatrix}
				\lambda_1I+N_1 & A_{12}& \cdots & A_{1s} \\ 
				& \lambda_2I+N_2 & \ddots &  \vdots\\ 
				&&\ddots & A_{s-1,s} \\
				&&&\lambda_sI+N_s \\ 
			\end{bmatrix}$. 取$A_1=\lambda_1I+N_1,A_2=\begin{bmatrix}
				 \lambda_2I+N_2 & \ddots &  \vdots\\ 
				&\ddots & A_{s-1,s} \\
				&&\lambda_sI+N_s \\ 
			\end{bmatrix},B=\begin{bmatrix}
			A_{12} & \cdots & A_{1s}
			\end{bmatrix}$, 由于$A_1,A_2$没有相同的特征值, 由上一小题, 存在$X_2=\begin{bmatrix}
			I & \wtd X_2\\ & I
			\end{bmatrix}^{-1}$, 使得$X_2X_1^{-1}AX_1X_2^{-1}=\begin{bmatrix}
				\lambda_1I+N_1 & & &  \\ 
				& \lambda_2I+N_2 & * &  *\\ 
				&&\ddots & A_{s-1,s} \\
				&&&\lambda_sI+N_s \\ 
			\end{bmatrix}$. 类似进行下去, 则有$X_s^{-1}\dots X_2^{-1} X_1^{-1}AX_1X_2\dots X_s=\begin{bmatrix}
				\lambda_1I+N_1 & & &  \\ 
				& \lambda_2I+N_2 & & \\ 
				&&\ddots &  \\
				&&&\lambda_sI+N_s \\ 
			\end{bmatrix}$. 令$X=X_1X_2\dots X_s$, 即得结论. 
            \comment{为何不直接使用Jordan标准型的结论？还是说还没教到Jordan标准型？}
		\item 记$X=\begin{bmatrix}
				X_1\\ X_2
				\end{bmatrix}$, 其中$X_2$是方阵. 记$\bs e_{1,t}=\begin{bmatrix}
				1 \\ 0 \\\vdots
			\end{bmatrix}_{t\times 1}$. 
		则方程
		\[
			J_m(0)X-XJ_n(0)=\bs e_1\bs b^T \Leftrightarrow
		\begin{bmatrix}
			J_{m-n}(0) & \bs e_{m-n}\bs e_{1,n}^T\\
			& J_n(0)
		\end{bmatrix}\begin{bmatrix}
			X_1\\ X_2
		\end{bmatrix}-\begin{bmatrix}
			X_1\\ X_2
		\end{bmatrix}J_n(0)=\begin{bmatrix}
		\bs e_{1,m-n}\bs b^T\\ \zero
		\end{bmatrix},
		\]
		可得
		\[
			J_{m-n}(0)X_1-\bs e_{m-n}\bs e_{1,n}^TX_2-X_1J_n(0)=\bs e_{1,m-n}\bs b^T,J_n(0)X_2-X_2J_n(0)=\zero.
		\]
		根据之前习题, $J_n(0)X_2-X_2J_n(0)=\zero$的解为
		\[
			X_2=t_1I+t_2J_n(0)+\dots+t_nJ_n(0)^{n-1}=\begin{bmatrix}
			t_1 & t_2 & \cdots & t_n\\ & t_1 & \ddots & \vdots\\
			&&\ddots &t_2\\ &&& t_1\\
		\end{bmatrix}=:\toep(\begin{bmatrix}
			t_1 & t_2 & \cdots & t_n\\
		\end{bmatrix})
		,
		\]
		由其第一行决定, 而第一行受前一方程约束. 
		前一方程得到$J_{m-n}(0)X_1-X_1J_n(0)=\bs e_{1,m-n}\bs b^T+\bs e_{m-n}\bs e_{1,n}^TX_2$. 
		如果$m-n=1$, 则有$\bs e_{1,n}^TX_2=-X_1J_n(0)-\bs b^T$, 
		任意选择$X_1$, 即$X$的第一行, 即可得到方程的解. 

		否则, 有
		\begin{align*}
			\bs b^T&=\bs e_{1,m-n}^TJ_{m-n}(0)X_1-\bs e_{1,m-n}^TX_1J_n(0)=\bs e_{2}^TX_1-\bs e_{1,m-n}^TX_1J_n(0),\\
			\bs 0^T&=\bs e_{p}^TJ_{m-n}(0)X_1-\bs e_{p}^TX_1J_n(0)=\bs e_{p+1}^TX_1-\bs e_{p}^TX_1J_n(0),(p=2,\dots,m-n-1)\\
			\bs e_{1,n}^TX_2&=\bs e_{m-n}^TJ_{m-n}(0)X_1-\bs e_{m-n}^TX_1J_n(0) =-\bs e_{m-n}^TX_1J_n(0).
		\end{align*}
		亦即, 
		\begin{align*}
			\bs e_{2}^TX_1&=\bs e_{1,m-n}^TX_1J_n(0)+\bs b^T,\\
			\bs e_{p}^TX_1&=\bs e_{p-1}^TX_1J_n(0),(p=3,\dots,m-n),\\
			\bs e_{1,n}^TX_2 &=-\bs e_{m-n}^TX_1J_n(0)
		.
		\end{align*}
		任意选择$X_1$的第一行, 即$X$的第一行, 即可逐行算出$X_1$以及$X_2$的第一行, 从而得到$X_2$, 从而完全得到$X$. 

		综上, 方程的解被其可任选的第一行所决定. 
		特别地, 第一行选作$\bs 0^T$, 则当$m-n=1$时, $X=\begin{bmatrix}
			\bs 0^T\\ \toep(\bs b^T)\\ 
		\end{bmatrix}$; 当$m-n>1$时, 
$X=\begin{bmatrix}
			I_{m-n}&\\& -I_n
		\end{bmatrix}\begin{bmatrix}
			\bs 0^T\\ \toep(\bs b^T)\\ \zero
		\end{bmatrix}$. 
			
			\item 取$A_1=J_{k_1}(0),A_2=J_{k_2}(0),\bs b=\bs a_2$, 由上一小题, 
				存在$\wtd X_2$, 使得
                \begin{align*}
                    X_2^{-1}AX_2=&\begin{bmatrix}
                        I & \wtd X_2 \\ & I \\ & &\ddots\\&&&I 
                    \end{bmatrix}
                    \begin{bmatrix}
            J_{k_1}(0) & \bs e_1\bs a_2^{\T} & \cdots & \bs  e_1\bs a_r^{\T}\\
            &J_{k_2}(0) \\
            && \ddots \\
            &&&J_{k_r}(0)\\
                    \end{bmatrix}\begin{bmatrix}
                        I & \wtd X_2 \\ & I \\ & &\ddots\\&&&I 
                    \end{bmatrix}^{-1}\\
                    =&\begin{bmatrix}
            J_{k_1}(0) & \zero & \cdots & \bs  e_1\bs a_r^{\T}\\
            &J_{k_2}(0) \\
            && \ddots \\
            &&&J_{k_r}(0)\\
                    \end{bmatrix}
                    .
                \end{align*}
				类似地, 
				存在$\wtd X_3$, 使得
                \begin{align*}
                    X_3^{-1}X_2^{-1}AX_2X_3=&\begin{bmatrix}
                        I & &\wtd X_3 \\ & I \\ & &I\\&&&\ddots 
                    \end{bmatrix}
                    \begin{bmatrix}
                        J_{k_1}(0) & \zero & \bs  e_1\bs a_3^{\T}&\cdots \\
            &J_{k_2}(0) \\
            && J_{k_3}(0)\\
            &&&\ddots \\
                    \end{bmatrix}\begin{bmatrix}
                        I & &\wtd X_3 \\ & I \\ & &I\\&&&\ddots 
                    \end{bmatrix}^{-1}\\
                    =&\begin{bmatrix}
                    J_{k_1}(0) & \zero & \zero & \cdots \\
            &J_{k_2}(0) \\
            && J_{k_3}(0)\\
            &&&\ddots \\
        \end{bmatrix}
        .
                \end{align*}
				进行下去, 有$X_r^{-1}\dotsm X_2^{-1}AX_2\dotsm X_r=\begin{bmatrix}
		J_{k_1}(0) \\
		&J_{k_2}(0)  \\
		&& \ddots \\
		&&&J_{k_r}(0)\\
	\end{bmatrix}$. 令$X=X_2\dotsm X_r$, 即得结论. 
	\end{enumerate}
\end{proof}

\begin{exercise}
设实对称矩阵$A=\begin{bmatrix}
a&1&1&-1 \\ 1&a&-1&1 \\ 1&-1&a&1 \\ -1&1&1&a \\
\end{bmatrix}$, 有一个单特征值$-3$, 求$a$的值和$A$的谱分解. 
\end{exercise}
\begin{proof}[解]
	关键是解出特征值, 并选择$a$使得$-3$是一个单特征值. 最朴素的做法是计算特征多项式
	\begin{align*}
		p_A(\lambda)=\det(\lambda I -A)
		&=\begin{vmatrix}
		\lambda-a & -1 & -1 & 1\\
		-1 & \lambda-a & 1 & -1\\
		-1 & 1 & \lambda-a & -1 \\
		 1& -1 & -1 & \lambda-a\\
		 \end{vmatrix}
		 \\&\xlongequal{\text{四行相加}}\begin{vmatrix}
		\lambda-a-1 & \lambda-a-1  & \lambda-a-1  & \lambda-a-1 \\
		-1 & \lambda-a & 1 & -1\\
		-1 & 1 & \lambda-a & -1 \\
		 1& -1 & -1 & \lambda-a\\
		 \end{vmatrix}
		 \\&=(\lambda-a-1)\begin{vmatrix}
		 1 & 1 & 1 & 1\\
		-1 & \lambda-a & 1 & -1\\
		-1 & 1 & \lambda-a & -1 \\
		 1& -1 & -1 & \lambda-a\\
		 \end{vmatrix}
		 \\&=(\lambda-a-1)\begin{vmatrix}
		 1 & 1 & 1 & 1\\
		 0 & \lambda-a+1 & 2 & 0\\
		 0 & 2 & \lambda-a+1 & 0 \\
		 2& 0 & 0 & \lambda-a+1\\
		 \end{vmatrix}
		 \\&\xlongequal{\text{对换二四行、二四列}}(\lambda-a-1)\begin{vmatrix}
		 1 & 1           & 1           & 1           \\
		 2 & \lambda-a+1 & 0           & 0           \\
		 0 & 0           & \lambda-a+1 & 2           \\
		 0 & 0           & 2           & \lambda-a+1 \\
		 \end{vmatrix}
		 \\&\xlongequal{\text{分块}}(\lambda-a-1)\begin{vmatrix}
		 1 & 1           \\
		 2 & \lambda-a+1 \\
		 \end{vmatrix}
		 \begin{vmatrix}
		 \lambda-a+1 & 2           \\
		 2           & \lambda-a+1 \\
	\end{vmatrix}
	\\&=(\lambda-a-1)(\lambda-a+1-2)\left((\lambda-a+1)^2-4\right)
	\\&=(\lambda-a-1)(\lambda-a+1-2)(\lambda-a+1+2)(\lambda-a+1-2)
	\\&=(\lambda-a-1)^3(\lambda-a+3)
	.
	\end{align*}可知特征值为$a+1,a+1,a+1,a-3$, 只有$a-3=-3,a+1\ne -3$, 即$a=0$时, $-3$是一个单特征值. 
	此时简单计算谱分解即得
	\[
		A=\begin{bmatrix}
		1 & 1& 1& 1\\
		-1 & 1 & 0 & 0\\
		-1 & 0 & 1 & 0\\
		1 & 0 & 0 & -1\\
	\end{bmatrix}\begin{bmatrix}
	-3 \\ & 1\\ && 1\\&&&1
	\end{bmatrix}\begin{bmatrix}
		1 & 1& 1& 1\\
		-1 & 1 & 0 & 0\\
		-1 & 0 & 1 & 0\\
		1 & 0 & 0 & -1\\
	\end{bmatrix}^{-1}
	.
\]	
一个稍微简化一点的做法是考虑$A$的特殊性, 得到$A=aI+B$, 其中$B=\begin{bmatrix}
1&1&1&-1 \\ 1&1&-1&1 \\ 1&-1&1&1 \\ -1&1&1&1 \\
\end{bmatrix}$. 已知条件转化为$B$有一个单特征值$-a-3$. 其他如上. 

第三种做法是利用定义, $A$有一个单特征值$-3$, 说明特征子空间$\nullspace(-3I-A)$的维数是$1$, 即$\rank(-3I-A)=3$. 于是
\begin{align*}
    -3I-A=&
	\begin{bmatrix}
			-3-a & -1 & -1 & 1\\
			-1 & -3-a & 1 & -1\\
			-1 & 1 & -3-a & -1 \\
			 1& -1 & -1 & -3-a\\
	\end{bmatrix}
	\to\begin{bmatrix}
			 1& -1 & -1 & -3-a\\
			0 & -4-a & 0 & -4-a\\
			0 & 0 & -4-a & -4-a \\
			0 & -4-a & -4-a & %1-(3+a)^2=
			(4+a)(-2-a)\\
	\end{bmatrix}\\
	&\quad\quad\to\begin{bmatrix}
			 1& -1 & -1 & -3-a\\
			0 & -4-a & 0 & -4-a\\
			0 & 0 & -4-a & -4-a \\
			0 & 0 & -4-a & %-(4+a)(2+a)+(4+a)=
			-(4+a)(1+a)\\
	\end{bmatrix}
	\to\begin{bmatrix}
			 1& -1 & -1 & -3-a\\
			0 & -4-a & 0 & -4-a\\
			0 & 0 & -4-a & -4-a \\
			0 & 0 & 0 & %-(4+a)(1+a)+(4+a)=
			-(4+a)a\\
	\end{bmatrix}
\end{align*}
容易看出：$a=0,\rank=3;a=-4,\rank=1$; $a$为其他值, 满秩. 
于是$a$只能是$0$. 
\end{proof}

\begin{exercise}
设三阶实对称矩阵$A$的各行元素之和均为$3$, $\bs a_1=\begin{bmatrix}
-1 \\ 2 \\ -1 \\
\end{bmatrix}, \bs a_2=\begin{bmatrix}
0 \\ -1 \\ 1 \\
\end{bmatrix}\in \nullspace(A)$, 求$A$及其谱分解. 
\end{exercise}
\begin{proof}[解]
	%记$A=\begin{bmatrix}
		%a_{11} & a_{12} & a_{13}\\
		%a_{21} & a_{22} & a_{23}\\
		%a_{31} & a_{32} & a_{33}\\
	%\end{bmatrix}$. 
	最朴素做法是直接利用已知条件, 然后求解$6$个未知数的方程组！
	好的做法是先动眼再动手. 稍微观察一下结构, 
	令$\bs 1=\begin{bmatrix}
		1\\1\\1
	\end{bmatrix}$, 则有$A\bs 1=3\cdot\bs 1, A\bs a_1=0,A\bs a_2=0$. 
	其中$\bs a_1,\bs a_2$线性无关. 
	于是$3$是一个几何重数至少为$1$的特征值, $0$是一个几何重数至少为$2$的特征值
	这就可以得到$A$的非对称谱分解
	\[
		A\begin{bmatrix}
			1 & -1 & 0\\
			1 & 2 & -1\\
			1 & -1 & 1\\
		\end{bmatrix}=\begin{bmatrix}
			1 & -1 & 0\\
			1 & 2 & -1\\
			1 & -1 & 1\\
		\end{bmatrix}\begin{bmatrix}
		3 \\ & 0 \\ &&0\\
		\end{bmatrix}.
	\]
		标准正交化后可得
		$A=\begin{bmatrix}
			\frac{1}{\sqrt{3}} & -\frac{1}{\sqrt{2}} & 0\\
			\frac{1}{\sqrt{3}} & \frac{1}{\sqrt{2}} & -\frac{1}{\sqrt{2}}\\
			\frac{1}{\sqrt{3}} & 0 & \frac{1}{\sqrt{2}}\\
		\end{bmatrix}\begin{bmatrix}
		3 \\ & 0 \\ &&0\\
		\end{bmatrix}\begin{bmatrix}
			\frac{1}{\sqrt{3}} & -\frac{1}{\sqrt{2}} & 0\\
			\frac{1}{\sqrt{3}} & \frac{1}{\sqrt{2}} & -\frac{1}{\sqrt{2}}\\
			\frac{1}{\sqrt{3}} & 0 & \frac{1}{\sqrt{2}}\\
		\end{bmatrix}^{\T}
		=\begin{bmatrix}
			1 & 1 & 1\\
			1 & 1 & 1\\
			1 & 1 & 1\\
		\end{bmatrix}
		$. 
\end{proof}

\begin{exercise}
求下列实对称矩阵$A$的谱分解：
\begin{enumerate}
\item $A$满足$A^3=\zero$. 
\item $A=a_1\bs x_1\bs x_1^{\T}+a_2\bs x_2\bs x_2^{\T}$, 其中$\bs x_1,\bs x_2$是$\mathbb{R}^2$的一组标准正交基, $a_1,a_2$为实数. 
\item $A=\begin{bmatrix}\zero&M\\M&\zero\end{bmatrix}$, 其中$M$是$n$阶对称矩阵, 有谱分解$M=Q\Lambda Q^{\T}$. 
\end{enumerate}
\end{exercise}
\begin{proof}[解]
	\begin{enumerate}
		\item 设$A=Q\Lambda Q^T$为谱分解, 其中$\Lambda$对角, $Q$正交. 则$A^3=Q\Lambda^3Q^T=\zero\Rightarrow \Lambda^3=\zero\Rightarrow \Lambda=\zero\Rightarrow A=\zero$. 
			\item $A=\begin{bmatrix}
					\bs x_1 & \bs x_2
			\end{bmatrix}\begin{bmatrix}
			a_1 & \\ & a_2\\
			\end{bmatrix}\begin{bmatrix}
					\bs x_1 & \bs x_2
			\end{bmatrix}^T$. 
		\item 
        \begin{align*}
            A=&\begin{bmatrix}
				\zero & Q\Lambda Q^T\\Q\Lambda Q^T & \zero
		\end{bmatrix}\\
        =&\begin{bmatrix}
			 Q & \\ & Q
		\end{bmatrix}\begin{bmatrix}
		\zero & \Lambda \\ \Lambda & \zero
		\end{bmatrix}\begin{bmatrix}
			 Q & \\ & Q
		\end{bmatrix}^T\\
        =&\begin{bmatrix}
			 Q & \\ & Q
		\end{bmatrix}\frac{1}{\sqrt{2}}\begin{bmatrix}
		I & I \\ I & -I
		\end{bmatrix}\begin{bmatrix}
		 \Lambda & \\ & -\Lambda  
		\end{bmatrix}\frac{1}{\sqrt{2}}\begin{bmatrix}
		I & I \\ I & -I
		\end{bmatrix}^T\begin{bmatrix}
			 Q & \\ & Q
		\end{bmatrix}^T.
        \end{align*}
	\end{enumerate}
\end{proof}

\begin{exercise}
设矩阵$A=\begin{bmatrix}2&b\\1&0\end{bmatrix}$. 
计算满足下列条件的$b$的取值范围：
\begin{enumerate}
\item $A$不可逆. 
\item $A$可以正交对角化. 
\item $A$不可对角化. 
\end{enumerate}
\end{exercise}
\begin{proof}[解]
	$A$不可逆, $A\ne0$, 说明$\rank(A)=1$. 第二列只能是第一列的倍数, 于是$b=0$. 
	另一种思路是$A$不可逆$\Leftrightarrow \det(A)=0$. 

	$A$可以正交对角化, 则$A$对称, 于是$b=1$. \note{如果在复数域上考虑, 则不能直接得到$A$对称, 只能得到$A$正规, 即$A^TA=AA^T$, 也可计算得到$b$. }

	$A$不可对角化, 则$A$有亏损特征值. 由于$A$最多两个不同特征值, 因此$A$只能是一个代数重数$2$几何重数$1$的特征值. 先计算$A$的特征多项式$p_A(\lambda)=\lambda^2-2\lambda-b$, 它有二重根时, $b=-1$, 此时重根为$1$. 然后再验证$\dim\nullspace(I-A)=1$, 或$\rank(I-A)=1$. 事实上, $\rank(I-A)=\rank\left(\begin{bmatrix}
		-1 & 1\\
		-1 & 1\\
	\end{bmatrix}\right)=1$. 
\end{proof}

\begin{exercise} \hard
设$A_1, A_2, \dots, A_m$是$m$个两两可交换的实对称矩阵, 证明它们可以同时正交对角化. 
\end{exercise}
\begin{proof}[证一]
	取$A_1$的一个特征值$\lambda_1$和特征子空间的一组基$\bs q_1,\dots,\bs q_t$, 将后者扩充成线性空间的一组基, 并按列组成矩阵$Q$, 则$Q^TA_1Q,\dots,Q^TA_mQ$两两可交换, 且$Q^TA_1Q=\begin{bmatrix}
		\lambda_1I & \\ &  A_{1,22}
	\end{bmatrix}$. 记$Q^TA_iQ=\begin{bmatrix}
	A_{i,11} & A_{i,21}^T \\ A_{i,21}^T & A_{i,22}
	\end{bmatrix}$, 由$A_1A_i=A_iA_1$知, $\begin{bmatrix}
	\lambda_1A_{i,11} & A_{i,21}^T\lambda_1 \\  A_{1,22}A_{i,21} &  A_{1,22}A_{i,22}
	\end{bmatrix}=\begin{bmatrix}
	\lambda_1A_{i,11} & A_{i,21}^T A_{1,22} \\ \lambda_1A_{i,21} & A_{i,22} A_{1,22}
	\end{bmatrix}$. 
	于是$ A_{1,22}A_{i,21}=\lambda_1A_{i,21}$. 
	而$\lambda_1$不是$ A_{1,22}$的特征值, 可得$ A_{1,22}-\lambda_1I$满秩, 因此$A_{i,21}=0$. 
	即$Q^TA_iQ=\begin{bmatrix}
	A_{i,11} &  \\  & A_{i,22}
	\end{bmatrix}$. 
	问题转化为$\lambda_1I, A_{2,11},\dots,A_{m,11}$和$ A_{1,22},A_{2,22},\dots,A_{m,22}$两个低阶问题, 对矩阵的阶利用数学归纳法即可证明. 
\end{proof}
\begin{proof}[证二]
	对$m$用数学归纳法. $m=1$时显然. 
	假设任意$m-1$个两两可交换的实对称矩阵可以同时对角化. 则存在$Q$使得$Q^TA_2Q=\begin{bmatrix}
		\lambda_{2,1}I \\ & \ddots\\ && \lambda_{2,t}I
	\end{bmatrix},\dots,Q^TA_mQ=\begin{bmatrix}
		\lambda_{m,1}I \\ & \ddots\\ && \lambda_{m,t}I
	\end{bmatrix}$, 其中分块满足\highlight[comment={这句话意义似乎不是很明确？如果我令$A_2,\cdots,A_m$全都等于单位阵呢？}]{至少有两个矩阵的对应对角块不同}. 
	记$Q^TA_1Q=\begin{bmatrix}
		A_{11} & \cdots & A_{t1}^T\\
		\vdots && \vdots \\
		A_{t1} & \cdots & A_{tt}
	\end{bmatrix}$, 利用交换性, 可知$(\lambda_{k,i}-\lambda_{k,j})A_{ij}=0,\forall k$. 当$i\ne j$时, 至少存在一个$k$, 使得$\lambda_{k,i}\ne \lambda_{k,j}$. （不然$Q^TA_kQ$的$(i,i)$块和$(j,j)$块全相同, 与之前假设矛盾. ）
	因此$A_{ij}=0$, 即$Q^TA_1Q=\begin{bmatrix}
		A_{11} \\ & \ddots \\ &&A_{tt}
	\end{bmatrix}$. 再对$A_{ii}$做谱分解$A_{ii}=P_{ii}^T\Lambda_i P_{ii}$, 其中$P_{ii}$正交, $\Lambda_i$对角. 
	记$P=\begin{bmatrix}
		P_{11} \\ &\ddots \\ && P_{tt}
	\end{bmatrix}$, 则$(QP)^TA_1(QP)=\begin{bmatrix}
		\Lambda_1 \\ &\ddots \\ && \Lambda_t
	\end{bmatrix},(QP)^TA_k(QP)=\begin{bmatrix}
		\lambda_{k,1}I \\ & \ddots\\ && \lambda_{k,t}I
	\end{bmatrix}$. 
\end{proof}

\clearpage
\listofchanges
\end{document}
